% ১০. Database Relation
\section{Database Relation}
রিলেশনাল ডেটাবেজে একাধিক টেবিল আলাদা রাখা হলেও প্রয়োজন অনুযায়ী তাদের মধ্যে সম্পর্ক তৈরি করা হয়। এই সম্পর্ক তৈরি হয় সাধারণত Primary key ও Foreign key-এর মাধ্যমে।

\subsection{Relation কী ও দরকারি ধরন}
\begin{definitionbox}{Primary Key ও Foreign Key}
\textbf{Primary Key} হলো এমন ফিল্ড যা প্রতিটি রেকর্ডকে অনন্যভাবে শনাক্ত করে। \textbf{Foreign Key} হলো এক টেবিলের এমন ফিল্ড, যা অন্য টেবিলের Primary key-এর সাথে সম্পর্ক তৈরি করে।
\end{definitionbox}

\begin{center}
\begin{tabular}{|p{0.24\textwidth}|p{0.45\textwidth}|p{0.20\textwidth}|}
\hline
রিলেশনের ধরন & ব্যাখ্যা & উদাহরণ \\
\hline
One-to-One & এক টেবিলের এক রেকর্ড অন্য টেবিলের এক রেকর্ডের সাথে যুক্ত & Person--Passport \\
\hline
One-to-Many & এক টেবিলের এক রেকর্ড অন্য টেবিলের একাধিক রেকর্ডের সাথে যুক্ত & Customer--Order \\
\hline
Many-to-Many & দুই পাশে একাধিক রেকর্ড যুক্ত; সাধারণত Junction table লাগে & Student--Course \\
\hline
\end{tabular}
\end{center}

\subsection{Junction table আলোচনা}
\begin{examplebox}{চিত্রের স্থানধারক}
Customer(CustomerID, Name) এবং Order(OrderID, CustomerID, Date) টেবিল দেখিয়ে CustomerID-এর মাধ্যমে One-to-Many relation আঁকা হবে।
\end{examplebox}

\begin{boardbox}{বোর্ড ফোকাস}
HSC পরীক্ষায় One-to-Many relation সবচেয়ে বেশি আসে। Parent table-এর Primary key child table-এ Foreign key হিসেবে ব্যবহার হয়।
\end{boardbox}
